\section{Root-poset cycle algorithms}
\label{app:root-poset-cycles}

We describe separately the absolute and non-Borel computations.  In both, a
source-leaf contraction removes a source of valency one together with its
unique target; isolated sources are discarded.  A surviving Jacobian with $r$
target rows and $c\ge r$ source columns is tested by enumerating the
$\binom cr$ source-column subsets.  Success means that one corresponding
determinant is a nonzero monomial in the simple-root variables.

\subsection{Absolute root posets}

The absolute routine applies the following finite test independently at each
height.

\begin{tcolorbox}
\begin{algorithmic}[1]
\Require An exceptional root system $\Phi$ and its Chevalley constants.
\State Construct the directed root poset and its bipartite height layers
$\fB(\Phi^+,\Delta,h)$.
\For{each height $h$ and each target selection $T$}
  \State Retain the full source layer and only edges ending in $T$.
  \State Repeatedly contract source leaves.
  \If{the graph contracts to the empty graph}
    \State Record a trivial certificate and continue.
  \EndIf
  \State Form the Jacobian
  $J_{\alpha,\beta}=N_{\beta,\alpha-\beta}x_{\alpha-\beta}$.
  \For{each set $C$ of $\#T$ source columns}
    \State Compute $d_C\gets\det J_C$.
    \If{$d_C$ is a nonzero monomial}
      \State Record $C$ and $d_C$ and continue with the next target selection.
    \EndIf
  \EndFor
  \State If no witness was found, report failure.
\EndFor
\end{algorithmic}
\end{tcolorbox}
\renewcommand{\figurename}{Algorithm}
\captionof{figure}{Absolute admissibility test}
\label{alg:absolute-cycles}

For the exceptional types the numerical summary is:
\[
\begin{array}{c|rr}
\toprule
\text{Type}&\text{target selections}&\text{failures}\\
\midrule
F_4&48&0\\
E_6&138&0\\
E_7&384&0\\
E_8&1196&0\\
\bottomrule
\end{array}
\]
The nontrivial determinants, up to orientation-dependent signs, are:
\[
\begin{array}{c|c|l}
\toprule
\text{Type}&\text{height}&\text{Jacobian minor}\\
\midrule
F_4&3&3x_{0001}x_{0010}x_{1000}\\
E_6&2&2x_{000010}x_{001000}x_{010000}\\
E_6&3&3x_{000001}x_{000010}x_{001000}x_{010000}x_{100000}\\
E_7&2&-2x_{0000100}x_{0010000}x_{0100000}\\
E_7&3&-3x_{0000010}x_{0000100}x_{0010000}x_{0100000}x_{1000000}\\
E_7&4&2x_{0000001}x_{0001000}x_{0010000}x_{0100000}\\
E_7&8&-2x_{0000001}x_{0000010}x_{0000100}x_{0010000}\\
E_8&2&-2x_{00001000}x_{00100000}x_{01000000}\\
E_8&3&-3x_{00000100}x_{00001000}x_{00100000}x_{01000000}x_{10000000}\\
E_8&4&2x_{00000010}x_{00010000}x_{00100000}x_{01000000}\\
E_8&5&-5x_{00000001}x_{00000010}x_{00000100}x_{00010000}x_{00100000}x_{01000000}x_{10000000}\\
E_8&8&-2x_{00000010}x_{00000100}x_{00001000}x_{00100000}\\
E_8&9&3x_{00000001}x_{00000100}x_{00001000}x_{00010000}x_{00100000}x_{10000000}\\
E_8&14&2x_{00000010}x_{00001000}x_{00010000}x_{01000000}\\
\bottomrule
\end{array}
\]

For $D_n$, all cycles are isomorphic to $\vec C_6$.  With the standard
ordering $\Delta_{D_n}=\{\delta_1,\ldots,\delta_n\}$, they are given by
\[
\begin{aligned}
\beta_1&=\sum_{i=n-h-1}^{n-2}\delta_i,&
\beta_2&=\sum_{i=n-h}^{n-2}\delta_i+\delta_{n-1},&
\beta_3&=\sum_{i=n-h}^{n-2}\delta_i+\delta_n,
\end{aligned}
\]
and their Jacobian coefficients are $2$ or $-2$.

\subsection{Non-Borel root posets}

Theorem~\ref{thm:twist-C6} already has the correct disjunction: a relative
cycle is either non-obstructing or its associated absolute graph is
admissible.  Thus the non-obstructing condition is applied before any absolute
Jacobian test.

\begin{tcolorbox}
\begin{algorithmic}[1]
\Require A Levi $\bM$, an elliptic Weyl element $w$, its relative root orbits,
and the non-obstructing singleton and pair relation.
\State Construct the $w$-relative root poset and its bipartite height layers.
\State Remove individually non-obstructing targets.
\For{each relative height layer}
  \State Contract source leaves.
  \For{each target subset $T$ containing no non-obstructing pair}
    \State Restrict to $T$ and contract source leaves again.
    \If{the graph contracts to the empty graph}
      \State Record the contraction and continue.
    \EndIf
    \State Replace every relative source and target by all absolute roots in
    its $w$-orbit.
    \State Add every absolute edge whose target-source difference is a
    positive root, and contract source leaves.
    \If{the absolute graph is nonempty}
      \State Apply Algorithm~\ref{alg:absolute-cycles}, testing combinations
      of absolute source columns.
    \EndIf
  \EndFor
\EndFor
\State Report failure if any surviving absolute graph has no monomial maximal
minor.
\end{algorithmic}
\end{tcolorbox}
\renewcommand{\figurename}{Algorithm}
\captionof{figure}{Non-Borel cycle test}
\label{alg:twisted-cycles}

The computation generates all exceptional Levi/Weyl configurations from the
Cartan matrices and covers:
\[
\begin{array}{c|rr}
\toprule
\text{Type}&\text{configurations}&\text{failures}\\
\midrule
F_4&19&0\\
E_6&68&0\\
E_7&145&0\\
E_8&304&0\\
\bottomrule
\end{array}
\]
Consequently all $F_4$ relative cycles are non-obstructing or contractible,
and every surviving E-type associated absolute graph is admissible.  The
reviewed data are available at
\url{https://github.com/mocham/AlgEG/tree/main/code/v4}.
